\subsection{Unit shortest arm and second shortest arm three}\label{proof:R2-C023}
For any finite residual list
$(1,3,c,s_2,t_2,\ldots,s_j,t_j)$,
$4\le c<s_2<t_2<\cdots<s_j<t_j$, fix the actual terminal
$F=qf_3h_c\prod_{i=2}^jh_{s_i}f_{t_i}$.
Orient every $L\ge1$ connector edge toward the receiving root and its
arm $a\ge3$ outward. Then
\begin{equation}\tag{R23.1.1}\label{eq:R23:claim}
\E[f_a(T^LF)^4]\le1-d/800.
\end{equation}
The $j=1$ row is the $(1,3,c)$ row of Section~\ref{proof:R2-C009}/$D_{\rm full}$.
The argument below treats every finite $j\ge2$ with the same
directions in both parameter ranges. Decorations may be restored.

Use the finite original receiving bounds, the exact covariance
\eqref{eq:R17:covariance}, and the Taylor formula
\eqref{eq:R17:taylor}.
$G_1$ has coordinates $(1,0)$; for $n\ge2$ the finite alternating
coordinates are in $[1-x,1]$. $\|G_n\|\le1$ for $n=1$ or $n\ge3$,
whereas $G_2$ has norm $\sqrt{1+x}<6/5$.
In particular $\operatorname{Cov}(C_a,\bar f_2)
\le-3d/16+\theta d^2$.
The full continuous cubic inequality \eqref{eq:R12:cubic}, with
$D=17/5$, $m\le33/32$, gives
\begin{equation}\tag{R23.2.5}\label{eq:R23:curvature}
C_a+Y+m\le17/5\ \Longrightarrow\
C_a[(Y+m)^2+2m^2]\le9,\qquad m^3\le9/8.
\end{equation}
Indeed $D\ge3m$, and its three terms are below $35/6$, $29/12$,
$11/20$, totaling $44/5<9$.
In the low range $0<d\le1/64$, each exponential payment below has
$C\ge A/2$, $A>B\ge0$, $Ad/2<1$. By
\eqref{eq:R14:exponential} the deficit exceeds $(A-B)d/3$.

The actual canonical first pair from Q5-C181, Section 4, has
\begin{equation}\tag{R23.3.1}\label{eq:R23:firstpair}
\begin{gathered}
f_3h_c\le e^{R_0-\Gamma_0},\quad \Gamma_0=d+\beta_c+\beta_{c-1}\ge35d/16,\quad R_0=\mathfrak r_0(B)u,\\
\mathfrak r_0=-\beta-z^3G_{c-3}(-z),\quad \beta=\beta_{c-1}.
\end{gathered}
\end{equation}
Write $G_{c-3}(-z)=U-zV$, with $(U,V)=(1,0)$ for $c=4$ and
$U,V\in[1-x,1]$ otherwise. Thus
$\mathfrak r_0=-\beta+x^2V+zxU$ and the actual signed mixed mean is
\begin{equation}\tag{R23.3.2}\label{eq:R23:mixedmean}
\E R_0+\E(uR_0)=-\beta d-\int x(U-xV)\,d\nu\le0,
\end{equation}
because $U-xV\ge1-2x\ge3/16$.
Every later actual pair has the envelope \eqref{eq:L96:long},
$h_sf_t\le e^{R_{s,t}-16d/5}$,
$R_{s,t}=\sum_{r=s+1}^tK^r\one$.
Let $S$ be their one complete finite interval mask, $R=R_0+S$.
Then
\[
F\le qe^{R-\Gamma},\qquad
\Gamma\ge431d/80+(16/5)(j-2)d.
\]
Since $s_2\ge5$, the first original power of $S$ is at least six.
The four residue classes of its finite zero--one mask give
\begin{equation}\tag{R23.3.4}\label{eq:R23:mask}
\|S\|_2^2\le\theta^5\frac{1+\theta}{(1-\theta^2)^2}d
=\frac{371293}{16634880}d<9d/400,\quad
\E S+\E(uS)\le\frac{2\delta_4}{1-\theta^2}<d/6.
\end{equation}
The positive mixed-mean contributions start with the paired original
$K^7,K^8$ terms $\delta_4$, with later same-sign indices spaced by two.
The geometric sums bound finite masks only. The row bound gives
$|S|<\sqrt d/4$ and
\[
\|R\|_2/\sqrt d\le d+(6/5)\theta^{3/2}+3/20
<777/1600<1/2,\qquad |R|\le d+2\sqrt d\le17/64<1/2.
\]
Therefore $e^R\le1+R+(3/5)R^2$. Expanding the actual $q=1+u$,
using the kernel of $T$ bounded by two and $q\le2$, gives
\begin{equation}\tag{R23.3.6}\label{eq:R23:firstrow}
T(qe^R)\le\bar f_2+BR+\delta,\qquad
\delta=\E R+\E(uR)+\sqrt d\|R\|_2+(12/5)\|R\|_2^2
\le19d/15<D:=4d/3.
\end{equation}
The $B(uR)$ cost is $\E|uR|$ and $T(qR^2)\le4\E R^2$.
The signed means $D_i=\E B^iR$ satisfy
\begin{equation}\tag{R23.3.7}\label{eq:R23:signed}
\begin{gathered}
|D_i|\le\tfrac12\theta^{(i-1)/2}d\ (i\ge1),\quad D_1\le d/2,\quad D_2\le d/3,\\
D_i\le\theta d/2\ (i\ge3).
\end{gathered}
\end{equation}
Finite propagation preserves the nonnegative upper bound
\begin{equation}\tag{R23.3.8}\label{eq:R23:propagation}
e^\Gamma T^LF\le Y_L=\bar f_{L+1}+B^LR+D\bar f_{L-1}
+\sum_{i=1}^{L-1}D_i\bar f_{L-1-i}\ge0.
\end{equation}

Two joint symbols retain the actual cross terms with $R_0$.
For $G_2+z\mathfrak r_0$, the coordinates are
$Q_e=1-x^2U$, $Q_o=1-\beta+x^2V$, both positive. Hence
\begin{equation}\tag{R23.4.1}\label{eq:R23:jointtwo}
|G_2+z\mathfrak r_0|^2
\le1+x-2x^2+4x^3+x^4-x^5+x^6
\le1+x+x^4\le1+\theta+\theta^4<36/25.
\end{equation}
Use $x<1/2$ for the two negative pairs; the last rational comparison
is $37578025<37748736$.
For $G_3+z^2\mathfrak r_0$ the coordinates are
$Q_e=1-x+\beta x-x^3V$, $Q_o=1-x^2U$, positive with lower
bounds $1-x-x^3$, $1-x^2$. Their square is bounded by
\begin{equation}\tag{R23.4.2}\label{eq:R23:jointthree}
[1-(1-\beta)x]^2+x
=1+x[-1+2\beta+(1-\beta)^2x]\le1,
\end{equation}
since $\beta\le1/64$, $x\le\theta$. Both symbols include $c=4$.

For $L=1$, $Y=\bar f_2+BR+D$ has actual mean
$m=1-d+D+D_1\le1+5d/6\le33/32$.
Its random vector is $[G_2+z\mathfrak r_0](B)u+BS$.
Using \eqref{eq:R23:jointtwo} and $\|BS\|_2<\sqrt d/10$ gives
$\operatorname{Var}Y\le169d/100$.
Since $\sigma(C_a)\le65\sqrt d/64$ and $\|BR\|_2<\sqrt d/3$,
\[
\operatorname{Cov}(C_a,Y)
\le-3d/16+\theta d^2+(65/64)(1/3)d
\le967d/6144<d/6.
\]
The pointwise row uses $|BR|\le\|R\|_2\le\sqrt d/2$, giving
$Y\le q+(3/2)\sqrt d+D$ and
$C_a+Y+m\le3+(27/10)\sqrt d+(13/6)d
\le3+713/1920<17/5$.
By \eqref{eq:R23:curvature}, the remainder is at most
$9(169/100)+4(9/8)(1/6)=399/25<16$.

For $L=2$, $Y=\bar f_3+B^2R+Dq+D_1$ has
$m\le1+d/6\le1+5d/6$.
Its random vector is $[G_3+z^2\mathfrak r_0](B)u+Du+B^2S$.
By \eqref{eq:R23:jointthree}, $D\le1/48$ and
$\|B^2S\|_2<\sqrt d/16$, its standard deviation is at most
$(13/12)\sqrt d$, and covariance at most
$(65/64)(13/12)d<9d/8$.
The correct three-step recurrence gives
$Y\le(1+D)q+(23/15)\sqrt d-d/2$ and
$C_a+Y+m\le3+(41/15)\sqrt d+(7/3)d
\le3+121/320<17/5$.
The remainder is at most
$9(169/144)+4(9/8)^2=125/8<16$.
For both rows the actual mean exponent is at least
$431/20+57/32-10/3=9599/480>19$, so
\begin{equation}\tag{R23.5.10}\label{eq:R23:shortpayment}
\E[f_a(T^LF)^4]\le e^{-19d}+16de^{-(431/20)d}.
\end{equation}
Its deficit exceeds $d>d/800$.

For $L=3,4$ raise coefficients only before nonnegative actual paths:
\[
\begin{aligned}
\overline Y_3&=\bar f_4+B^3R+D\bar f_2+(d/2)q+d/3,\\
\overline Y_4&=\bar f_5+B^4R+D\bar f_3+(d/2)\bar f_2+(d/3)q+d/4.
\end{aligned}
\]
They dominate \eqref{eq:R23:propagation} and are nonnegative.
The correct bare recurrences and vector polynomials are those of
Modules R2-C014 and R2-C017, with means at most $1-83d/32$ and
$1-11d/4$, and norms $\|G_4-d\|\le1$,
$\|G_5-d(2+z)\|\le1-2d$.
At $L=3$, retain $\E B^3R\le\theta d/2$ to get
$m\le1-[83/32-4/3-1/2-1/3-\theta/2]d=1-43d/192\le1$.
The whole-vector norm is at most
$[1+(21/10)d+2/15]\sqrt d<(6/5)\sqrt d$, including the true
$G_2$ cost. The row bound is
$\overline Y_3\le(1+5d/6)q+(5/4)\sqrt d-2d/3$,
so $C_a+\overline Y_3+m\le3+d+(49/20)\sqrt d<17/5$.
At $L=4$, $\E B^4R<(2/15)d$ gives
$m\le1-[11/4-4/3-1/2-1/3-1/4-2/15]d=1-d/5$.
The whole-vector norm is at most
$[1+(34/15)d+1/12]\sqrt d<(6/5)\sqrt d$.
The actual row bound is
$\overline Y_4\le(1+d/6)q+(6/5)\sqrt d-\tau+d/4$;
the scalar remainder is negative, so
$C_a+\overline Y_4+m\le3+d/3+(12/5)\sqrt d<17/5$.
Both rows have variance at most $36d/25$, covariance at most $5d/4$,
and remainder at most $9(36/25)+5=449/25<18$.
The mean exponent is at least $431/20+57/32=3733/160>23$, giving
\begin{equation}\tag{R23.6.8}\label{eq:R23:middlepayment}
\E[f_a(T^LF)^4]\le e^{-23d}+18de^{-(431/20)d}.
\end{equation}
Its deficit exceeds $5d/3>d/800$.

For all finite $L\ge5$, put $n=L+1\ge6$. The term
$D\bar f_{L-1}$ has index two, where $\beta_2=0$, and is paid
explicitly at cost $2D=8d/3$.
For $r\ge3$, $\gamma_r=\beta_r-D_{r-2}$ satisfies
$\gamma_3\ge d/2$, $\gamma_4\ge2d/3$, and
$\gamma_r\ge(\kappa-\theta/2)d=25d/64\ge h=3d/8$ for $r\ge5$.
The additional actual mean index satisfies
$\beta_{n+1}-D_L\ge h$ as well.
Retaining the last four actual paths gives the nonnegative bound
\[
e^\Gamma T^LF\le\widehat Y
=1-h(4-d)+[G_n-h(3+2z+z^2)](B)u+B^LR+8d/3.
\]
The exact mean estimate is
$m\le1-5h+3hd+8d/3
=1+(19/24)d+(9/8)d^2\le1+5d/6\le33/32$.
The full finite shifted-symbol certificates
\eqref{eq:R18:shiftfour}--\eqref{eq:R18:shifttwo} apply.
Since $\|B^LR\|_2\le(1/2)\theta^{L/2}\sqrt d<\sqrt d/16$,
variance is at most $289d/256$ and covariance at most
$(1105/1024)d<(11/10)d$.
The pointwise form is $(1-3h)q+BV-h(1-d)+8d/3$, where
$V=[G_L-h(2+z)](B)u+B^{L-1}R$ and
$\|V\|_2\le(1-2h+1/12)\sqrt d\le(13/12)\sqrt d$.
As $dp+3hq\ge2d$,
\[
C_a+\widehat Y+m\le3+(9/8)d+(3/8)d^2+(137/60)\sqrt d
\le3+(6/5)d+(12/5)\sqrt d<17/5.
\]
The remainder is at most
$9(289/256)+4(9/8)(11/10)=19341/1280<16$.
Thus \eqref{eq:R23:shortpayment} covers every $L\ge5$ too,
completing the entire low range.

For $d\ge1/64$, retain the actual first three arms.
Every later pair has $s\ge5$ and its original interval satisfies
$|R_{s,t}|\le(6/5)\theta^{(s-1)/2}\sqrt d
<\sqrt d/5\le8d/5$.
Its loss is at least $16d/5$, so $h_sf_t\le e^{-8d/5}$.
Consequently, with unchanged directions,
\[
F\le e^{-(8/5)(j-1)d}qf_3h_c,\qquad
T^LF\le e^{-(8/5)(j-1)d}T^L(qf_3h_c).
\]
The exact Module Section~\ref{proof:R2-C009}/$D_{\rm full}$ row, originating in Q5-C175,
gives $\E[f_a\{T^L(qf_3h_c)\}^4]\le1-d/800$ for every
$a\ge3,c\ge4,L\ge1$.
The full target is therefore at most
$e^{-(32/5)(j-1)d}(1-d/800)\le1-d/800$.
This joins the low range at $d=1/64$; $j=1$ uses the same directory
row directly. At $d=0$ original leaf integration gives equality,
without $K=0$. Positivity restores original equal pairs and full
reversal reverses the receiving arm, first three arms, all subsequent
pairs, connectors, and decorations together.
